% from ega2-indexes-and-contents-fr.tex
\oldpage[II]{208}
\markboth{A. GROTHENDIECK --- Chap. II}{A. GROTHENDIECK --- Chap. II}
$\beta,\beta_{\mathcal{F}}$ (homomorphismes canoniques) : \hyperref[II.3.3.4-fr]{3.3.4}.

$G(\vphi),\Proj(\vphi)$
($\vphi$ homomorphisme de $\mathcal{O}_Y$-Algèbres graduées) : \hyperref[II.3.5.1-fr]{3.5.1}.

$\Proj(u)$ ($u$ $g$-morphisme d'une $\mathcal{O}_Y$-Algèbre graduée
dans une $\mathcal{O}_{Y'}$-Algèbre graduée) : \hyperref[II.3.5.6-fr]{3.5.6}.

$r_{\mathcal{L},\psi}$ ($\mathcal{L}$ $\mathcal{O}_X$-Module inversible) : \hyperref[II.3.7.1-fr]{3.7.1}.

$\mathbf{P}(\mathcal{E}),\mathbf{P}(E),\mathbf{P}_Y^n,\mathbf{P}_A^n$ : \hyperref[II.4.1.1-fr]{4.1.1}.

$\mathbf{P}(u)$ ($u$ homomorphisme surjectif de $\mathcal{O}_Y$-Modules) : \hyperref[II.4.1.2-fr]{4.1.2}.

$\Hyp_Y(X,\mathcal{E})$ ($\mathcal{E}$ $\mathcal{O}_Y$-Module quasi-cohérent) : \hyperref[II.4.2.5-fr]{4.2.5}.

$\varsigma$ (morphisme de Segre) : \hyperref[II.4.3.1-fr]{4.3.1}.

$\mathcal{F}(n)$ ($\mathcal{F}$ $\mathcal{O}_X$-Module, $X$ muni d'un
$\mathcal{O}_X$-Module inversible) : \hyperref[II.4.5.1-fr]{4.5.1}.

$\varepsilon$ (automorphisme identique) : \hyperref[II.4.5.1-fr]{4.5.1}.

$P(u,T),\sigma_i(u)$ ($u$ endomorphisme de module libre de base finie) : \hyperref[II.6.4.1-fr]{6.4.1}.

$P(u,T),\sigma_i(u),\det u$
($u$ endomorphisme de module de type fini sur un anneau intègre ou un
anneau noethérien réduit) : \hyperref[II.6.4.2-fr]{6.4.2} et \hyperref[II.6.4.7-fr]{6.4.7}.

$\sigma_i(u),\det u$ ($u$ endomorphisme de $\mathcal{A}$-Module localement libre) : \hyperref[II.6.4.8-fr]{6.4.8}.

$\sigma_i,\det$ (homomorphismes de faisceaux) : \hyperref[II.6.4.8-fr]{6.4.8}, \hyperref[II.6.4.9-fr]{6.4.9} et \hyperref[II.6.4.10-fr]{6.4.10}.

$\sigma_i(u),\det u$
($u$ endomorphisme de $\mathcal{O}_X$-Module de type fini sur un préschéma
localement intègre ou un préschéma localement noethérien réduit) :
\hyperref[II.6.4.9-fr]{6.4.9} et \hyperref[II.6.4.10-fr]{6.4.10}.

$N_{\mathcal{B}/\mathcal{A}}(f)$
($f$ section de la $\mathcal{B}$-Algèbre $\mathcal{A}$) : \hyperref[II.6.5.1-fr]{6.5.1}.

$N_{\mathcal{B}/\mathcal{A}}(\mathcal{L}')$
($\mathcal{L}'$ $\mathcal{B}$-Module inversible) : \hyperref[II.6.5.2-fr]{6.5.2}.

$N_{\mathcal{B}/\mathcal{A}}(h')$
($h'$ homomorphisme de $\mathcal{B}$-modules inversibles) : \hyperref[II.6.5.3-fr]{6.5.3}.

$N_{X'/X}(\mathcal{L}')$
($X'$ fini au-dessus de $X$, $\mathcal{L}'$ $\mathcal{O}_{X'}$-Module inversible) : \hyperref[II.6.5.5-fr]{6.5.5}.

$N_{X'/X}(h')$
($h'$ homomorphisme de $\mathcal{O}_{X'}$-Modules inversibles) : \hyperref[II.6.5.5-fr]{6.5.5}.

$\mathcal{S}_X$ ($X=\Proj(\mathcal{S})$, $\mathcal{S}$ somme d'Idéaux de
$\mathcal{O}_Y$) : \hyperref[II.8.1.5-fr]{8.1.5}.

$S^\geq,S^\leq,S_f^\geq,S_f^\leq,
M^\geq,M^\leq,M_f^\geq,M_f^\leq$
($S$ anneau gradué, $M$ $S$-module gradué,
$f$ élément homogène de $S_+$) : \hyperref[II.8.2.1-fr]{8.2.1}.

$\widehat S$ ($S$ anneau gradué) : \hyperref[II.8.2.2-fr]{8.2.2}.

$\widehat M$ ($M$ $S$-module gradué) : \hyperref[II.8.2.4-fr]{8.2.4}.

$S_{[n]},S^\natural,f^\natural$
($S$ anneau gradué, $f$ homogène dans $S$) : \hyperref[II.8.2.6-fr]{8.2.6}.

$M_{[n]},M^\natural$ ($M$ $S$-module gradué) : \hyperref[II.8.2.8-fr]{8.2.8}.

$\widehat{\mathcal{S}}$ ($\mathcal{S}$ $\mathcal{O}_Y$-Algèbre graduée) : \hyperref[II.8.3.1-fr]{3.8.1}.

$G_m$ (schéma en groupes) : \hyperref[II.8.3.9-fr]{8.3.9}.

$\mathcal{S}_X$ ($X=\Proj(\mathcal{S})$, $\mathcal{S}$ $\mathcal{O}_Y$-Algèbre graduée) : \hyperref[II.8.6.1-fr]{8.6.1}.

$\mathcal{S}_X,C_X,\widehat C_X,E_X,\widehat E_X$
($X=\Proj(\mathcal{S})$, $\mathcal{S}$ $\mathcal{O}_Y$-Algèbre graduée
avec $\mathcal{S}_0=\mathcal{O}_Y$) : \hyperref[II.8.6.1-fr]{8.6.1}.

$\mathcal{S}_{[n]},\mathcal{S}^{\natural}$
($\mathcal{S}$ $\mathcal{O}_Y$-Algèbre graduée) : \hyperref[II.8.7.2-fr]{8.7.2}.

$\egaShProj_0(\mathcal{M}),\egaShProj(\mathcal{M}),\mathcal{M}_X,
\widehat{\mathcal{M}},\mathcal{M}^{\square}$
($\mathcal{M}$ $\mathcal{S}$-Module gradué) : \hyperref[II.8.12.1-fr]{8.12.1}.

$\mathcal{M}_X^\geq$
($X=\Proj(\mathcal{S})$, $\mathcal{M}$ $\mathcal{S}$-Module gradué) : \hyperref[II.8.12.5-fr]{8.12.5}.

$\mathcal{M}_{[n]},\mathcal{M}^{\natural}$
($\mathcal{M}$ $\mathcal{S}$-Module gradué) : \hyperref[II.8.12.9-fr]{8.12.9}.

$(\widetilde{\mathcal{N}})^-$
($\mathcal{N}$ sous-$\mathcal{S}$-Module d'un $\mathcal{S}$-Module gradué) : \hyperref[II.8.13.1-fr]{8.13.1}.

$\overline{\mathcal{N}}$
($\mathcal{N}$ sous-$\mathcal{S}$-Module d'un $\mathcal{S}$-Module gradué) : \hyperref[II.8.13.2-fr]{8.13.2}.

$\Gamma_*(\mathcal{F})$
($X=\Proj(\mathcal{S})$, $\mathcal{F}$ $\mathcal{S}_X$-Module gradué) : \hyperref[II.8.14.5-fr]{8.14.5}.

$\mathcal{S}_{X'},\Gamma'_*(\mathcal{F}'),\egaShProj'(\mathcal{M}),
\alpha',\beta'$
($X'$ ouvert de $X=\Proj(\mathcal{S})$, $\mathcal{F}'$ $\mathcal{S}_{X'}$-Module gradué,
$\mathcal{M}$ $\mathcal{S}$-Module gradué) : \hyperref[II.8.14.6-fr]{8.14.6}.

\endgroup

\clearpage
